
\section{Herramientas y entorno de desarrollo}
\subsection{Lenguajes de programación}
El repositorio oficial de SerenityOS utiliza múltiples lenguajes. De acuerdo con las estadísticas de lenguajes publicadas por GitHub (endpoint \url{https://api.github.com/repos/SerenityOS/serenity/languages}), el volumen principal del código está en C++ y, en menor proporción, en otros lenguajes usados para subsistemas, tooling y automatización. En los valores reportados (en bytes) se observan, entre otros:

\begin{itemize}
		\item C++: \texttt{46966564}
		\item JavaScript: \texttt{2387389}
		\item HTML: \texttt{1775193}
		\item Shell: \texttt{442848}
		\item CMake: \texttt{440499}
		\item C: \texttt{435237}
		\item Python: \texttt{225053}
		\item Assembly: \texttt{59192}
\end{itemize}

Para construir herramientas de host, la guía oficial de compilación indica un compilador en el host con soporte de C++23; versiones probadas incluyen GCC 13 y Clang 17 a 19.

\subsection{Compilador cruzado}
Para esta parte del informe, el entorno de compilación utilizado es un sistema Arch Linux de 64 bits (x86\_64). La guía oficial de SerenityOS define que el proceso de build incluye un compilador cruzado (toolchain) que se construye automáticamente la primera vez que se ejecuta el flujo de compilación.

Dependencias en Arch Linux:

\begin{verbatim}
sudo pacman -S --needed \
	base-devel cmake curl mpfr libmpc gmp e2fsprogs ninja qemu-desktop \
	qemu-system-aarch64 ccache rsync unzip

# Opcional:
# - fuse2fs (imágenes sin root)
# - clang llvm llvm-libs (para compilar con Clang)
\end{verbatim}

Clonado del repositorio y navegación:
\begin{verbatim}
git clone https://github.com/SerenityOS/serenity.git
cd serenity
\end{verbatim}

Construcción del toolchain y compilación/ejecución (driver \texttt{Meta/serenity.sh}):
\begin{verbatim}
# Compila todo y ejecuta en QEMU.
Meta/serenity.sh run

# Compila sin arrancar la VM.
Meta/serenity.sh build
\end{verbatim}

En la primera ejecución de \texttt{Meta/serenity.sh run}, la guía indica que el script (i) descarga archivos de base requeridos desde Internet y (ii) construye la toolchain cruzada; estos pasos se realizan una sola vez y el resto de builds se ejecutan más rápido.

Estructura de salida y rutas relevantes descritas por la guía:
\begin{itemize}
		\item Directorio de instalación de artefactos: \texttt{Build/<architecture>/Root}
		\item Arquitecturas soportadas por el flujo: \texttt{x86\_64}, \texttt{aarch64}, \texttt{riscv64}
\end{itemize}

En el caso \texttt{x86\_64}, la configuración recomendada para herramientas de análisis estático (p.
ej., clangd) utiliza como referencia el \texttt{compile\_commands.json} bajo \texttt{Build/x86\_64/compile\_commands.json} y el compilador cruzado ubicado bajo \texttt{Toolchain/Local/x86\_64/bin/}, con ejecutables como \texttt{x86\_64-pc-serenity-g++}.

\subsection{Emulación y virtualización}
La guía de compilación indica que \texttt{Meta/serenity.sh run} inicia SerenityOS usando QEMU, además de construir el sistema e instalar los artefactos bajo \texttt{Build/<architecture>/Root}. En la primera ejecución, el tiempo extra está asociado a la descarga de bases requeridas y a la construcción de la toolchain cruzada. En ejecuciones posteriores (con toolchain ya preparada) el flujo típico es: compilar incrementalmente y lanzar directamente la VM.

Durante \texttt{Meta/serenity.sh run} el script puede solicitar credenciales de administrador mediante \texttt{sudo} para operaciones del host relacionadas con la ejecución (según configuración del sistema). En el flujo normal, al finalizar, se abre una ventana de QEMU y el sistema operativo se emula/arranca automáticamente.

Requisito de versión: QEMU 6.2 o superior. Si el sistema no dispone de esa versión, la guía contempla compilar una versión recomendada mediante \texttt{Toolchain/BuildQemu.sh}.

Dependencias adicionales para compilar QEMU en Debian/Ubuntu (según la guía):
\begin{verbatim}
sudo apt install libgtk-3-dev libpixman-1-dev libsdl2-dev libslirp-dev \
	libspice-server-dev
\end{verbatim}

Opcional: \texttt{fuse2fs} para construir imágenes sin privilegios de administrador. La guía documenta que, si aparece el error \texttt{fusermount: failed to open /etc/mtab}, se crea el enlace \texttt{/etc/mtab} hacia \texttt{/proc/self/mounts}.

\subsection{Control de versiones}
Se utiliza Git como sistema de control de versiones para el informe y para el seguimiento del código del sistema base. El repositorio oficial de SerenityOS se encuentra en GitHub: \texttt{https://github.com/SerenityOS/serenity}.

\subsection{Editor y entorno de desarrollo}
Para desarrollo de código, SerenityOS no impone un IDE único: es posible trabajar con editores como Vim, Emacs o entornos gráficos.

Configuración y herramientas de VS Code recomendadas por la documentación de SerenityOS:
\begin{itemize}
		\item Extensiones sugeridas: \texttt{clangd}, \texttt{GitLens}, y soporte para Jakt (si se usa).
		\item Configuración de workspace: ajustes en \texttt{.vscode/settings.json} para excluir directorios generados (\texttt{Build/}, \texttt{Toolchain/Local/}, etc.) y parámetros de clangd.
		\item Base para IntelliSense/diagnóstico: \texttt{compile\_commands.json} generado en \texttt{Build/<arch>/compile\_commands.json}.
\end{itemize}

Formateo y estilo (según la guía \textit{Serenity C++ coding style}): el estilo de bajo nivel se deriva del archivo \texttt{.clang-format} en la raíz del repositorio, y se indica usar \texttt{clang-format} versión 20 o superior.

Para el documento de este informe, se utiliza Visual Studio Code con LaTeX Workshop para compilación y vista previa.